\begin{pastexam}[label={px:nccu-113-dsa-04}]{政大資科所 113 學年度 第 4 題}
（單選；不必說明理由）

(4\%) There are various designs for min heaps (e.g., binary heap, binomial heap, and Fibonacci heap), and
they all support ExtractMin, DecreaseKey, and Insert.

Which of the following statements about min heaps is correct:
\begin{enumerate}[label=\Alph*.]
\item DecreaseKey always takes $\Theta(\log n)$ time in a binary heap.
\item Insert always takes $O(1)$ time in a binary heap.
\item It is possible to design a min heap such that both DecreaseKey and Insert always take $O(1)$ time.
\item It is possible to design a min heap such that both ExtractMin and Insert always take $O(1)$ time.
\end{enumerate}
\end{pastexam}

\begin{pastsoln}
答案是 C。

\begin{itemize}
\item[A.] 錯。 二元堆積的 DecreaseKey 是把鍵改小之後往上浮（資結書虛擬碼~8.1、8.4）。 改完之後如果還是不小於
父節點（例如改的是根，或只改小一點點），一次比較就停，$O(1)$。 $\Theta(\log n)$ 是最壞情況，不是「每次都」。
\item[B.] 錯。 二元堆積的 Insert 放在最後再往上浮，插入比所有鍵都小的數時要浮到根，$\Theta(\log n)$
（資結書第~8.1~節）。
\item[C.] 對。 費氏堆積的 Insert 和 DecreaseKey 是\textbf{平攤} $O(1)$（資結書第~8.6~節）。 題目說
「always」，要的是每次都 $O(1)$；這也做得到：Brodal（1996）的堆積在最壞情況下 Insert、DecreaseKey 都是
$O(1)$，ExtractMin 是 $O(\log n)$。 所以「可以設計出」這樣的最小堆積，敘述成立。
\item[D.] 錯。 如果 Insert 和 ExtractMin 都是 $O(1)$，把 $n$ 個數全部插入、再取出 $n$ 次，就在 $O(n)$ 時間內
排好序了。 這個堆積只靠比較鍵的大小，排序要 $\Omega(n \lg n)$ 次比較（資結書定理~15.1），矛盾。 平攤 $O(1)$
也不行，同樣的論證適用於總時間。
\end{itemize}
\end{pastsoln}
